% problem_id: S001-lemma-polarized-diagonal-separated-fp
% source_id: S001 (Stacks Project)
% source_file: moduli.tex
% statement_locator: moduli.tex lines 1765--1769
% proof_locator: moduli.tex lines 1771--1892
% env: lemma
% id_note: label 在本来源内唯一
% pairing: tight (verified)
% status: complete (statement + author proof verbatim + reason + locators)
%
\begin{lemma}
\label{lemma-polarized-diagonal-separated-fp}
The diagonal of $\Polarizedstack$ is separated
and of finite presentation.
\end{lemma}

\begin{proof}
Recall that $\Polarizedstack$ is a limit preserving algebraic stack, see
Quot, Lemma \ref{quot-lemma-polarized-limits}.
By Limits of Stacks, Lemma \ref{stacks-limits-lemma-limit-preserving-diagonal}
this implies that
$\Delta : \Polarizedstack \to \Polarizedstack \times \Polarizedstack$
is limit preserving. Hence $\Delta$ is locally of finite presentation
by Limits of Stacks, Proposition
\ref{stacks-limits-proposition-characterize-locally-finite-presentation}.

\medskip\noindent
Let us prove that $\Delta$ is separated. To see this, it suffices to show
that given an affine scheme $U$ and two objects
$\upsilon = (Y, \mathcal{N})$ and $\chi = (X, \mathcal{L})$
of $\Polarizedstack$ over $U$, the algebraic
space
$$
\mathit{Isom}_{\Polarizedstack}(\upsilon, \chi)
$$
is separated. The rule which to an isomorphism $\upsilon_T \to \chi_T$
assigns the underlying isomorphism $Y_T \to X_T$ defines a morphism
$$
\mathit{Isom}_{\Polarizedstack}(\upsilon, \chi)
\longrightarrow
\mathit{Isom}_U(Y, X)
$$
Since we have seen in Lemmas \ref{lemma-Mor-s-lfp} and
\ref{lemma-Isom-in-Mor} that the target is
a separated algebraic space, it suffices to prove that this morphism
is separated. Given an isomorphism $f : Y_T \to X_T$
over some scheme $T/U$, then clearly
$$
\mathit{Isom}_{\Polarizedstack}(\upsilon, \chi)
\times_{\mathit{Isom}_U(Y, X), [f]} T
=
\mathit{Isom}(\mathcal{N}_T, f^*\mathcal{L}_T)
$$
Here $[f] : T \to \mathit{Isom}_U(Y, X)$ indicates the $T$-valued
point corresponding to $f$ and
$\mathit{Isom}(\mathcal{N}_T, f^*\mathcal{L}_T)$
is the algebraic space discussed in Section \ref{section-hom-isom}.
Since this algebraic space is affine over $U$, the claim implies
$\Delta$ is separated.

\medskip\noindent
To finish the proof we show that $\Delta$ is quasi-compact. Since
$\Delta$ is representable by algebraic spaces, it suffice to check
the base change of $\Delta$ by a surjective smooth morphism
$U \to \Polarizedstack \times \Polarizedstack$ is quasi-compact
(see for example Properties of Stacks, Lemma
\ref{stacks-properties-lemma-check-property-covering}).
We can assume $U = \coprod U_i$ is a disjoint union of affine opens.
Since $\Polarizedstack$ is limit preserving (see above), we
see that $\Polarizedstack \to \Spec(\mathbf{Z})$ is locally of
finite presentation, hence $U_i \to \Spec(\mathbf{Z})$ is
locally of finite presentation
(Limits of Stacks, Proposition
\ref{stacks-limits-proposition-characterize-locally-finite-presentation}
and Morphisms of Stacks, Lemmas
\ref{stacks-morphisms-lemma-composition-finite-presentation} and
\ref{stacks-morphisms-lemma-smooth-locally-finite-presentation}).
In particular, $U_i$ is Noetherian affine. This reduces us to the
case discussed in the next paragraph.

\medskip\noindent
In this paragraph, given a Noetherian affine scheme $U$ and two objects
$\upsilon = (Y, \mathcal{N})$ and $\chi = (X, \mathcal{L})$
of $\Polarizedstack$ over $U$, we show the algebraic space
$$
\mathit{Isom}_{\Polarizedstack}(\upsilon, \chi)
$$
is quasi-compact. Since the connected components of $U$ are open and closed
we may replace $U$ by these. Thus we may and do assume $U$ is connected.
Let $u \in U$ be a point. Let $P$ be the Hilbert polynomial
$n \mapsto \chi(Y_u, \mathcal{N}_u^{\otimes n})$, see
Varieties, Lemma \ref{varieties-lemma-numerical-polynomial-from-euler}.
Since $U$ is connected and since
the functions $u \mapsto \chi(Y_u, \mathcal{N}_u^{\otimes n})$
are locally constant (see 
Derived Categories of Schemes, Lemma
\ref{perfect-lemma-chi-locally-constant-geometric})
we see that we get the same Hilbert polynomial in every point of $U$.
Set
$\mathcal{M} = \text{pr}_1^*\mathcal{N}
\otimes_{\mathcal{O}_{Y \times_U X}} \text{pr}_2^*\mathcal{L}$
on $Y \times_U X$. Given
$(f, \varphi) \in \mathit{Isom}_{\Polarizedstack}(\upsilon, \chi)(T)$
for some scheme $T$ over $U$ then for every $t \in T$ we have
$$
\chi(Y_t, (\text{id} \times f)^*\mathcal{M}^{\otimes n}) =
\chi(Y_t,
\mathcal{N}_t^{\otimes n} \otimes_{\mathcal{O}_{Y_t}}
f_t^*\mathcal{L}_t^{\otimes n}) =
\chi(Y_t, \mathcal{N}_t^{\otimes 2n}) = P(2n)
$$
where in the middle equality we use the isomorphism
$\varphi : f^*\mathcal{L}_T \to \mathcal{N}_T$.
Setting $P'(t) = P(2t)$ we find that the morphism
$$
\mathit{Isom}_{\Polarizedstack}(\upsilon, \chi)
\longrightarrow
\mathit{Isom}_U(Y, X)
$$
(see earlier) has image contained in the intersection
$$
\mathit{Isom}_U(Y, X) \cap \mathit{Mor}^{P', \mathcal{M}}_U(Y, X)
$$
The intersection is an intersection of open subspaces of
$\mathit{Mor}_U(Y, X)$ (see Lemma \ref{lemma-Isom-in-Mor} and
Remark \ref{remark-Mor-numerical}).
Now $\mathit{Mor}^{P', \mathcal{M}}_U(Y, X)$
is a Noetherian algebraic space as it is of finite
presentation over $U$ by Lemma \ref{lemma-Mor-qc-over-base}.
Thus the intersection
is a Noetherian algebraic space too. Since the morphism
$$
\mathit{Isom}_{\Polarizedstack}(\upsilon, \chi)
\longrightarrow
\mathit{Isom}_U(Y, X) \cap \mathit{Mor}^{P', \mathcal{M}}_U(Y, X)
$$
is affine (see above) we conclude.
\end{proof}
